\section{Data and method}
\label{sec:method}

\subsection{Data}
\label{sec:data}

We used four keystroke datasets released for research, from three sources (Table~\ref{tab:data}). Two are editions of an
introductory programming (CS1) course at a US university, in which students wrote Python in a browser-based editor and
every edit was logged in the same format \cite{edwards2023cs1}; the 2019 and 2021 editions share the course, the
editor and the recording. CS1 records edit events rather than key presses; an insertion of one character became a
character press and a deletion a deletion press (S2). KLiCKe \cite{tian2025klicke} is a corpus of argumentative essays
written in 2022 by adults in the United States, recruited online, with English as a first or second language; each
wrote one essay on one of four prompts adapted from the SAT, with a 30-minute timer and the submit button disabled until
30 minutes had passed. Trained raters scored each essay on a six-point holistic rubric (the mean of at least two raters, in
steps of 0.5; weighted Cohen's $\kappa$ 0.60 before and 0.76 after discussing scores two or more points apart)
\cite{tian2025klicke}. The task, a timed argumentative essay scored on a holistic rubric, has the format of writing
assessment in education, although the writers were not sampled from a course. CoAuthor is a dataset of creative and argumentative writing with suggestions from a language model
\cite{lee2022coauthor}; it is writing outside an educational setting and enters as a contrast with a different
recording, not as an educational sample. The method needs two fields per key press: its time and whether it inserted
or deleted a character.

A writer is a student (CS1) or a participant (KLiCKe, CoAuthor). We excluded the units of writers with less than 10
minutes of observation (the time within the observation intervals, not the time on the KLiCKe timer) and, in the count analyses, units with
less than 5 minutes of exposure. KLiCKe is the only dataset with an outcome that we used: in the 2019 CS1 release the
grade table was withdrawn by the data authors as erroneous, and the 2021 grades were not used because only 31 writers
were on the seen side. Holdouts set aside in earlier work have no confirmatory role here (Table~\ref{tab:data}).

\begin{table}[t]
\centering\footnotesize
\caption{The four datasets. Units are assignments (CS1), essays (KLiCKe) and sessions (CoAuthor). Onsets are the
stall onsets under the baseline definition.}
\label{tab:data}
\begin{tabularx}{\textwidth}{@{}l X r r r X X@{}}
\toprule
Dataset & Writers and setting; one event & Writers & Units & Onsets & Outcome used & Access and holdout \\
\midrule
CS1 2019 & students of a CS1 course writing Python; one logged edit & 459 & 1,469 & 33,710 & none (grade table
withdrawn) & from the data authors on request, aggregate values only; 146 held-out writers opened earlier, included \\
CS1 2021 & same course, editor and recording & 31 & 177 & 6,217 & none (31 writers) & as CS1 2019; 13 held-out
writers not opened \\
KLiCKe & adults writing one argumentative essay; one key press & 4,908 & 4,908 & 33,003 & holistic score (1--6)
& public, free use with citation \cite{inputlog2026community}; 1,498 held-out writers opened earlier, included \\
CoAuthor & adults writing with model suggestions; one key press & 41 & 1,155 & 1,518 & none & public, released for
research; 18 held-out writers (288 sessions) not opened \\
\bottomrule
\end{tabularx}
\end{table}

\paragraph{Relation to the companion study.} A companion study \cite{miura2026runup} tested, on a holdout, whether an
estimator's signals distinguish the run-up to a stall from a short pause. The two studies share the baseline
definition, the conversion of CS1, the observation intervals written by that estimator and the opened holdouts, so
baseline rates and part of $\varphi$ can coincide with values reported there. Everything else here is new: the 89
definitions and the pause family, the intervals, the score association and its adjustments, the overdispersion within
writers and the length decomposition (S8).

\subsection{The baseline definition and the detector}
\label{sec:detector}

A stall is detected from three ingredients (Figure~\ref{fig:detector}). Let \emph{condition A} hold at time $t$ when,
in the window $(t - W, t]$, the number of character presses minus the number of deletion presses is less than one. When
condition A rises from false to true at time $s$ and stays true until $t_0 = s + \theta$, an onset is raised at $t_0$;
$\theta$ is the value at $s$ and is held until $t_0$, condition A is evaluated only from $W$ after the start of an
observation interval, $t_0$ must fall within the interval, and a new onset requires condition A to become false and
rise again.
The stall lasts from $s$ to the first character press after $t_0$ in the same observation interval (the
resumption); its length is
measured from $t_0$, so it can be short if typing resumes soon after the onset. Under the characters-only rule an onset
requires no character press for at least $W + \theta$, more than 15~s in our grid; under the net rule, deletions can make
condition A hold while characters are still being typed, so the label is a stop in net character production and
retyping counts as a stop. The detector does not reach pauses of a few seconds, which the pause family below covers. The \emph{baseline definition} is
the default of the detector's reference implementation: $W = 30$~s, $\theta$ the 0.99 nearest-rank quantile of the
writer's inter-keystroke intervals so far in the unit (defined once 100 intervals exist), the net rule for condition A,
and the default observation intervals. It is the definition of the companion study's appendix and was fixed before any
KLiCKe score was read. The detector reproduced the onsets of the reference implementation in every unit of the four
datasets (S2).

The default observation intervals are the spans in which the companion study's estimator, re-run on the key stream,
wrote its per-second rows; they cover nearly all of a KLiCKe essay, with holes of a few seconds (S2). Because they
depend on the estimator, the grid also contains two segmentations defined by key times alone, and every result is
reported for them. A writer-quantile $\theta$ is undefined before the unit's 100th interval, so no stall can start in
that time (3.2--10.3\% of the exposure), whereas a fixed $\theta$ can; we repeated the comparisons with all
definitions restricted to the time after the 100th interval (S2). Condition A is evaluated only from $W$ after the
start of an observation interval, so no definition can detect a stall in the first $W$ seconds of an interval
(0.5--10.9\% of the exposure across datasets, segmentations and windows); we also repeated the comparisons with each
interval starting 60~s after its start and not before the 100th interval (S2).

\begin{figure}[t]
\centering
\begin{tikzpicture}[x=0.11cm,y=0.5cm,font=\footnotesize]
  \draw[->] (0,0) -- (120,0) node[right] {time};
  \foreach \x in {2,5,7,10,13,15,18,21,24,27} \draw (\x,0) -- (\x,0.45);
  \foreach \x in {96,99,103,106,110} \draw (\x,0) -- (\x,0.45);
  \node[above] at (14,0.5) {character presses};
  \draw[fill=figband,draw=none] (27,-0.35) rectangle (57,-0.95);
  \node at (42,-0.65) {$W$};
  \draw[thick] (57,0.3) -- (57,-2.2) node[below] {$s$};
  \draw[thick] (81,0.3) -- (81,-2.2) node[below] {$t_0$};
  \draw[<->] (57,-1.55) -- node[above] {$\theta$} (81,-1.55);
  \draw[thick] (96,0.3) -- (96,-2.2) node[below] {resumption};
  \draw[decorate,decoration={brace,mirror,amplitude=4pt}] (57,-3.3) -- node[below=4pt] {stall (length measured
  from $t_0$)} (96,-3.3);
\end{tikzpicture}
\caption{The stall detector. Condition A (fewer than one net character press in the window $W$, shaded) becomes true at
$s$; if it stays true for $\theta$, an onset is raised at $t_0 = s + \theta$. The stall ends at the next character
press.}
\label{fig:detector}
\end{figure}

\subsection{The four axes and the 89 definitions}
\label{sec:grid}

We varied four choices that every such detector has to make, each over levels drawn from prior practice or set around
the default (Table~\ref{tab:grid}). The full crossing gives 90 definitions; with the baseline removed, 89 remain, and
all proportions below are out of 89 and describe this grid, not a sample of possible definitions. The two families of
$\theta$ are different constructs: a quantile of the writer's own intervals marks a pause that is long \emph{for this
writer}, a fixed number of seconds marks a pause that is long in absolute terms. We report them together because both
are in use, and separately where they part (53 definitions with a writer quantile, 36 with fixed seconds).

\begin{table}[t]
\centering\footnotesize
\caption{The four axes of the stall detector. The baseline level is marked with an asterisk. The grid of default and
precedented levels keeps the default window and condition, which have no use outside the detector, and the levels with
precedent outside it: writer-relative thresholds as a kind, the 10-s threshold and the three segmentations (12
definitions, including the baseline). A fixed $\theta$ is a hold after condition A rises, so a 10-s $\theta$ is not the
same construct as a 10-s pause threshold.}
\label{tab:grid}
\begin{tabularx}{\textwidth}{@{}l l X@{}}
\toprule
Axis & Levels & Source of the levels \\
\midrule
Window $W$ & 15, 30$^\ast$, 60~s & 30~s is the default; 15 and 60~s are half and double the default; we found no
use of a window length outside the detector \\
Threshold $\theta$ & writer quantile 0.95, 0.99$^\ast$, 0.995; fixed 10, 30~s & writer-relative thresholds are in use
(an ex-Gaussian bound per writer \cite{chukharev2014pauses}, ten times the within-word median \cite{guo2020scenario}),
though not at these quantiles; fixed thresholds in writing span about 0.25--10~s \cite{torrance2024methods}, so 10~s
has precedent and 30~s, set equal to the default window, has none \\
Condition A & net$^\ast$ (characters $-$ deletions $< 1$); characters only ($< 1$) & whether deletions cancel character
presses; neither rule has a use outside the detector \\
Observation & default intervals$^\ast$; split at key-press gaps of 5~min; of 30~min & 5~min of inactivity as
time-on-task \cite{hart2023time}; 30~min as a session timeout \cite{halfaker2015session} \\
\bottomrule
\end{tabularx}
\end{table}

\paragraph{The pause family.} The second detector counts a pause whenever two consecutive key presses in the same
observation interval are more than a threshold apart. We crossed two kinds of interval (between any two key presses,
or between two character presses), three fixed thresholds (2, 5 and 10~s, within the range used in writing
\cite{torrance2024methods}) and the three segmentations, for 18 definitions. We also computed 18 definitions with the
writer's 0.95, 0.99 and 0.995 quantile over the unit as the threshold; such a threshold labels a fixed share of each
writer's intervals, so its rate is close to that share times the writer's rate of key presses, and we report these as
a control for the key-press rate rather than as pause labels (S3). The pause family is not a check of the stall
results but the same diagnostic applied to a second detector.

\subsection{What we measured}
\label{sec:measures}

\paragraph{Rates and ratios.} The rate is the number of stalls or pauses per observed hour, pooled over a dataset's
units. For each definition we compared the rank order of the four datasets with the baseline order, and the ratio of
the rates of each pair of datasets with the baseline ratio.

\paragraph{The essay-score association.} In KLiCKe, the Spearman correlation of a writer's rate with the
holistic score. Because the sample is large, weak correlations have intervals that exclude zero, so we also describe magnitudes,
using $|\rho| < 0.1$ as negligible and 0.3 as moderate \cite{cohen1988power}.

\paragraph{Adjustments.} Two partial Spearman correlations of score and rate: given typing speed, the median of a
writer's intervals shorter than 2~s, measured six ways (S3); and given the essay's length in characters and the length
of observation. Essay length is part of what raters judge, so the length-adjusted value is the association with the
part of the rate that does not overlap with length, not an effect.

\paragraph{Overdispersion within writers.} For sub-writer units (CS1 work blocks and assignments; CoAuthor sessions),
the Pearson dispersion index $\varphi = \sum (n - \hat e)^2 / \hat e \,/\, (k - p)$, where the expected count $\hat e$ is
the writer's own rate times the unit's exposure and $p$ the number of writers. $\varphi = 1$ is what counting at a
constant rate within a writer would give; we read $\varphi$ as the departure from a constant rate, with 1.5 as a
working threshold, not as a choice between count models, because the window of the detector, dependence between
stalls and the mean structure all affect it (S4).

\paragraph{Components of stall length.} A linear mixed model of log stall length with random effects of writer,
assignment, day and work block (CS1; days are runs of blocks separated by gaps of 6~h or more) or of writer and
session (CoAuthor), fitted by REML; we report the share of each level and of the residual, the stall-to-stall
variation. Stalls that did not resume within their interval are excluded (S5). We do not report the analogous
decomposition of counts, whose estimates did not recover simulated components (S6).

\paragraph{Intervals.} All intervals are 95\% percentile intervals from 2,000 resamples of writers, with each
writer's labels as detected. Two-stage intervals that also redraw the writer-quantile $\theta$ are a
sensitivity analysis (S2). The interval of a
definition shows sampling variation; the spread across definitions shows definitional variation, and we set the two
side by side \cite{robitzsch2022multiverse}. Definitions share axes and are not independent, so we read counts as
proportions of the grid and do not test the grid as a whole \cite{simonsohn2020specification}.

\subsection{Status of the analyses}
\label{sec:status}

Every computation ran under conditions committed in the author's repository before it ran; these are not third-party
registrations. S1 gives the order of the computations, including those added after seeing results,
and the status of each conclusion. The speed covariate was chosen twice after seeing results, so the speed-adjusted values are exploratory. Three points bear
on the claims. The criterion for a reversal of a rate ratio (25\% of the
definitions for a pair) was set after the number of reversed point values per pair had been seen, and the pair committed
to be reported first, the two CS1 editions, did not reach it. The criterion for a reversal of the score association
required only one definition with an interval on the other side of zero, which is lenient, so we report the full
distribution; requiring the point value to lie there as well, a rule added later, changed no count. Sizes, the
reading by classes of interval and the grid of default and precedented levels were set after seeing the results.

The author used Claude (Anthropic) to search the literature, to write the analysis scripts under these conditions and
to check the numbers in the text against the result files; every number in the text was compared with the result files.
